% ---------------------------------------------------------------------------
% Section 4 -- The Framework
%
% Budget: 1,400-1,700 words (WRITING-GUIDE.md Sec.5). Largest prose section.
%
% Narrative contract inherited from the Introduction:
%   * tension 1 "the inventory does not exist"  -> paid off in Sec.4.2
%   * tension 3 "is the model just guessing?"   -> set up in Sec.4.3/4.4,
%     paid off by the 55/55 result in the Evaluation, NOT here.
%
% Every number below appears verbatim in GROUND-TRUTH.md. Citations are sparse
% by design (WRITING-GUIDE.md Sec.6: method sections cite only what grounds a
% design decision).
%
% Terminology: "resource" throughout, matching Sec.1 and Sec.3 (Table
% tab:observables writes the resource r). The policy's inventory of resources
% is "the register".
%
% NOTE for the next build pass: this section resolves \ref{sec:threat},
% \ref{eq:decision} and \ref{sec:eval}, so 03-05 must all be \input in
% main.tex before the cross-references typeset.
% ---------------------------------------------------------------------------
\section{The Framework}
\label{sec:framework}

\framework{} computes the score $\rho(c)$ of \S\ref{sec:threat} in two layers.
A static layer builds a baseline for every (tool, resource) pairing a server
exposes; a dynamic layer refines it with the concrete arguments,
the session, and the calling agent's own history. The output is an ordinal
band.

\subsection{Overview}

The static layer reads exactly two documents: the tool catalog the server
advertises over \texttt{tools/list}, and the section of the organization's
policy that covers this server. Neither is a \framework{}-specific artifact:
the server already publishes the catalog, and a security team already maintains
the policy. The policy supplies a register of the resources behind the server
and the tools that touch them, named classes, and
the rules by which the organization's own staff recognise them; it states no
severity number anywhere. The organization's per-resource sensitivity table
never reaches the scorer. Here it is an evaluation instrument only --- the
artifact a deployment would not have, which is what lets it serve as an oracle.
\S\ref{sec:eval} separates the real captures from what was written for this
study.

Three primitives, each on a 1--5 scale, are estimated and multiplied:
\begin{equation}
\mathrm{score} = \mathrm{sensitivity} \times \mathrm{blast} \times
                 \mathrm{impact}, \qquad \mathrm{score} \in [1, 125].
\label{eq:score}
\end{equation}
The three ask deliberately different questions: impact is about the verb, what
one call does; sensitivity about the noun, what the call touches; blast radius
about the quantifier, how far the consequence travels. Multiplying rather than
averaging keeps a low factor influential: a destructive tool aimed at a
worthless resource cannot reach the top of the range, since only a call severe
on every axis gets there.

A scan issues four \emph{kinds} of prompt: stage~0 infers the server's domain
once, fixing the vocabulary for what follows; stage~1 assigns tool impact per
tool, rules first and the model only on abstention; stage~2 assigns resource
sensitivity per register row; stage~3 assigns blast radius per (tool, resource)
cell. The last stage dominates a scan's prompt count --- 26 tools over 20
resources is 520 cells. Everything afterwards is deterministic. An assembly pass reconciles
twin tools (a bulk variant of a singular tool, or a deprecated alias of a
canonical one, is tied to its twin so neither can be priced below the other)
and applies a blast floor; a banding function then maps the product onto
$\{\textrm{low}, \textrm{medium}, \textrm{high}, \textrm{critical}\}$. Banding
is not a raw threshold: it encodes explicit security floors --- any
irreversible operation is at least medium --- derived by measuring where an LLM
reviewer systematically disagreed with raw thresholds. Determinism buys exactly
one thing, and we claim only that: band boundaries, twin reconciliation and the
floor do not drift between runs. The residual run-to-run variance sits in the
LLM stages, and \S\ref{sec:eval} quantifies it for blast radius, the noisiest
of them.

\subsection{The Static Layer}
\label{sec:static}

The static layer runs once per server, before any agent connects. It scores
every (tool, resource) pairing the catalog exposes, producing the baseline
that the dynamic layer later refines. Its three primitives are estimated
independently and combined by Eq.~\ref{eq:score}.

\subsubsection{Deriving Sensitivity from Policy}
\label{sec:sensitivity}

The input a risk model most wants is the one an organization is least able to
publish. A per-resource sensitivity inventory is expensive to build, sensitive
in itself --- filenames can be protected health information --- and stale the
moment a new resource appears. What organizations do publish is a data
classification policy.

\framework{} therefore derives sensitivity by \emph{classify-then-map},
splitting the labour so neither party does the other's job. The organization
states each class's consequence of loss in its own words ---
statutory breach, patient injury, privilege waiver, recoverable embarrassment
--- and writes no numbers. \framework{} owns the 1--5 scale, maps that
adverse-impact language onto it, and classifies each register row into one of
the organization's classes, using the row's description and the policy's own
recognition and aggregation rules. Classify-then-map follows the security
categorization procedure of FIPS~199 and
SP~800-60~\cite{nist2004fips199,nist2008sp80060} --- impact language first,
category second --- applied to a document the organization already maintains
rather than to a labelled inventory. A rule, unlike an inventory row, still
classifies a resource created after the policy was written --- a property of
the construction, not a result we measure.

\subsubsection{Tool Impact: A Deterministic Ladder}
\label{sec:impact}

Tool impact asks one question of a tool's own declaration: is this operation a
read, a write, or a removal? Table~\ref{tab:ladder} is the resulting ladder,
matched by a vocabulary of 281 patterns --- 16 for \emph{none}, 52 metadata, 46
read, 122 write, 45 remove --- plus 119 ambiguous single words matched against
the tool name only, 10 generic read verbs, and 9 parameter detectors. That last
scoping matters: matched across a description instead, the ambiguous words make
a listing tool look destructive.

\begin{table}[t]
\caption{The tool-impact ladder. One question --- read, write, or removal ---
decides the tier.}
\label{tab:ladder}
\small
\begin{tabular}{@{}cp{0.80\columnwidth}@{}}
\toprule
Tier & Operation \\
\midrule
1 & \textbf{none} --- the server talks about itself: ping, health, version,
    whoami \\
2 & \textbf{metadata} --- names, ids, counts, sizes, timestamps, permissions,
    schema, listings; consumption state \\
3 & \textbf{content read or limited write} --- append a line, add a comment,
    set one named field \\
4 & \textbf{ordinary write} --- the caller supplies what the item says \\
5 & \textbf{removal or execution} --- delete, wipe, drop, purge, truncate;
    execute code; move money \\
\bottomrule
\end{tabular}
\end{table}

The tier 3/4 boundary --- limited write against ordinary write --- is not
invented here: two published standards draw it in the same place. HTTP
separates \texttt{PATCH}, a partial modification leaving unmentioned fields
untouched~\cite{dusseault2010patch}, from \texttt{PUT}, where the client sends
the complete representation~\cite{fielding2022http}; CVSS~v4.0 separates
integrity impact \texttt{VI:L}, ``the amount of modification is limited'', from
\texttt{VI:H}, ``a total loss of integrity''~\cite{first2023cvss4}. Writing a
file is a 4; writing one sentence into it is a 3.

Breadth is deliberately absent from the ladder: how many items a call reaches
is coverage, which blast radius scores, so the rules read no breadth vocabulary
(\emph{all}, \emph{every}, \emph{bulk}, \emph{batch}), no glob and no recursion
switch. Parameter detectors are recorded as capability
flags but never move a tier: what a caller \emph{could} pass is not what a call
\emph{does} pass, which the dynamic layer sees and the static layer cannot. MCP tool annotations are
recorded the same way, as evidence and never as a
bound~\cite{mcp2026spec} --- an annotation is a component's self-assessment
rather than an independent measurement, and no component sets the ceiling on
its own risk score.

\subsubsection{The Rules-to-Model Hand-off}
\label{sec:handoff}

The rules report a confidence with every tier: \textbf{0.35} when no tier verb
matched, \textbf{0.8} when a verb fired, and \textbf{0.95} when a verb fired and
an MCP annotation corroborated it. The model is consulted only below
\textbf{0.5}. An annotation therefore raises confidence, and can suppress the
model's review, but never sets or caps the tier itself.

The model is asked where the rules do not know --- a write whose declaration
never says whether it is bounded --- not where two rules disagree, which the
precedence order settles. What the
arrangement buys is attributability rather than call volume: stages~2 and~3
already spend a prompt per cell, so the tool axis is not where a scan's prompt
volume lives. Every tier the rules decide traces to a named pattern, and a per-tool
field records which scorer decided, so no tier is unattributable. On the
catalogs we evaluate the fallback never fired; \S\ref{sec:eval} reports that
result and what it stops us claiming.

\subsubsection{Blast Radius}
\label{sec:blast}

Blast radius is scored per (tool, resource) pairing: how far one call's
consequences propagate, which subjects and systems are different afterwards.
The rubric numbers that propagation. (1)~Almost nobody depends on what was
touched; (2)~a few subjects; (3)~a group, the people or systems
attached to what was touched, and the normal case; (4)~the whole population
inside the resource boundary, the subjects the register says one resource
holds; (5)~beyond that boundary. Tier~5 requires a route the organization's own
register sanctions: the resource is a hub other systems
authenticate against, what the call returns is usable on its own elsewhere, one
call reaches the resource's entire population, or the resource is destroyed
with nothing left to restore. Counting items would break on the obvious case:
reading one credential file touches one item and compromises every system that
credential opens.

One deterministic pass then acts on the number: a floor ties blast radius to
the sensitivity of the resource and to irreversible operations. That breaks the
primitives' independence on purpose, in one direction only --- sensitivity can
raise blast, never lower it --- because a call reaching a crown-jewel resource
is not a pinpoint consequence whatever its verb. An earlier version also
applied a roof; it was removed, because a roof can only
ever under-score. Blast radius remains the noisiest of the three stages, which
\S\ref{sec:eval} quantifies rather than hides.

\subsection{The Dynamic Layer}
\label{sec:dynamic}

The static table answers how bad a call would be if it happened; it cannot see
the agent, the session, or the arguments actually supplied. The dynamic layer
answers how abnormal this call is, from this agent, in this session, given its
actual input, using four signals ordered cheap to expensive.

\textbf{(0) Embedding likelihood.} A hashing vectorizer over the tool name,
argument tokens and structural flags, reduced with TruncatedSVD fit on the
server's observed history; novelty is the mean distance to the 5 nearest
historical calls, calibrated against the history's own quantiles.
\textbf{(1) Behavioral baseline.} Deviation from the calling agent's own
history: a tool it has never called, a first touch of a high-sensitivity
resource, a burst above its normal rate.
\textbf{(2) Session sequence risk.} Cumulative sensitivity within a session,
plus a flag for a sensitive read followed by an external write in a window ---
a pattern no per-call view can see.
\textbf{(3) LLM judge.} Decodes and normalizes arguments (base64, hex, shell
expansion) and classifies intent --- the only signal that uses a model, and it
is optional.

Three of the four are LLM-free, which keeps the layer's compute cost
negligible: fitting takes under 0.5\,s per server and scoring roughly 0.05\,ms
per call. The judge is wired but not exercised here.

Two compositions coexist, entering at different points. The embedding
likelihood fills a likelihood slot,
\begin{equation}
\mathrm{final\_risk}(c) = \mathrm{score}(c) \times \mathrm{likelihood}(c),
\qquad \mathrm{likelihood} \in [0.1, 1.0],
\label{eq:dynamic}
\end{equation}
so that demonstrated-normal behaviour discounts risk rather than erasing it.
The 0.1 floor makes that safe: a cell the scan called critical keeps a tenth of
its risk however routine the call looks, and a call with no benign history
keeps full static risk at likelihood 1.0, so an agent never observed is never
silently discounted. The three band signals $s_{1}, s_{2}, s_{3}$ --- signals
(1)--(3) above --- then fuse by escalation on top of the banded result, each
able only to \emph{raise} it; none can argue a dangerous capability down.
Banding first and escalating second gives the quantity \S\ref{sec:threat}
asked for:
\begin{equation}
\rho(c) = \mathrm{escalate}\bigl(\mathrm{band}(\mathrm{final\_risk}(c)),\,
          s_{1}(c),\, s_{2}(c),\, s_{3}(c)\bigr),
\label{eq:rho}
\end{equation}
an ordinal band over the same four labels.

\subsection{Enforcement}
\label{sec:enforcement}

Enforcement instantiates Eq.~\ref{eq:decision}: $\rho(c)$ is the band just
defined and $\theta_{r} \le \theta_{d}$ are the organization's own, so each is
an operating point rather than a redesign. \S\ref{sec:eval} measures the
flag/no-flag boundary at both a \emph{high} and a \emph{critical} setting, and
reports the benign fall-out each costs. The intermediate \textsf{review}
outcome --- holding a call for human verification rather than refusing it ---
belongs to the design and is not evaluated here.

The decision is auditable in the organization's own terms. A refusal decomposes
into the three primitives of Eq.~\ref{eq:score}, the policy class the resource
was mapped to (\S\ref{sec:sensitivity}), the operation tier and the scorer that
produced it (\S\ref{sec:handoff}), and the likelihood that scaled the result.
Two of the three primitives are model-assigned, so this is not a model-free
audit trail; what the decomposition gives an operator is a named class and a
named operation to dispute, in place of a single opaque number.
